% GENERATED FILE. Edit canonical lessons, then run npm run generate:books.
% canonical-source-hash: fnv1a64:c5b75b37360ec84f
% canonical-lessons: SA-C177-pracchadayati, SA-C177-vitarati, SA-C177-sangrhnati, SA-C177-visvasiti, SA-C177-lajjate

\chapter{Conceals, Distributes, Collects, Trusts, Is ashamed}
\label{ch:sa-conceals-distributes-collects-trusts-is-ashamed}
\hlchaptermodality{177}

\begin{chapteropening}
\textbf{By the end of this chapter:} \emph{I can say conceals, distributes, collects, trusts and is ashamed.}

Complete the last lesson of chapter 177: I can say conceals, distributes, collects, trusts and is ashamed.
\end{chapteropening}

\section[\texorpdfstring{\sk{प्रच्छादयति} (pracchādayati)}{प्रच्छादयति (pracchādayati)}]{\sk{प्रच्छादयति} (pracchādayati) — conceals, covers}

\label{lesson:SA-C177-pracchadayati}



\begin{quote}
Before the new one: say the Sanskrit for cleans, then the Sanskrit for fills.
\end{quote}

\subsection*{You'll want to know: \sk{प्रच्छादयति}}
\textbf{\sk{प्रच्छादयति}} — \emph{pracchādayati} — \textquotedblleft{}conceals, covers\textquotedblright{}.

\begin{grammarlens}[title={What you've built}]
One more word for this chapter.
\end{grammarlens}

\subsection*{Your turn}
\begin{itemize}
\raggedright
  \item \emph{Say it:} \emph{pracchādayati}
  \item \emph{Say it:} \emph{pracchādayati}, once more
  \item \emph{Say it:} say \emph{pracchādayati}
  \item \emph{Recall:} say the Sanskrit for cleans, then the Sanskrit for fills, then say \emph{pracchādayati} again
\end{itemize}

\begin{tcolorbox}[breakable,colback=teal!4,colframe=teal!35!black,title={Before you move on}]
What does \sk{प्रच्छादयति} mean? (\textquotedblleft{}Conceals, covers\textquotedblright{}.) Say it once more.
\end{tcolorbox}
\section[\texorpdfstring{\sk{वितरति} (vitarati)}{वितरति (vitarati)}]{\sk{वितरति} (vitarati) — distributes, hands out}

\label{lesson:SA-C177-vitarati}



\begin{quote}
Before the new one: say the Sanskrit for fills, then the Sanskrit for conceals.
\end{quote}

\subsection*{You'll want to know: \sk{वितरति}}
\textbf{\sk{वितरति}} — \emph{vitarati} — \textquotedblleft{}distributes, hands out\textquotedblright{}.

\begin{grammarlens}[title={What you've built}]
One more word for this chapter.
\end{grammarlens}

\subsection*{Your turn}
\begin{itemize}
\raggedright
  \item \emph{Say it:} \emph{vitarati}
  \item \emph{Say it:} \emph{vitarati}, once more
  \item \emph{Say it:} say \emph{vitarati}
  \item \emph{Recall:} say the Sanskrit for fills, then the Sanskrit for conceals, then say \emph{vitarati} again
\end{itemize}

\begin{tcolorbox}[breakable,colback=teal!4,colframe=teal!35!black,title={Before you move on}]
What does \sk{वितरति} mean? (\textquotedblleft{}Distributes, hands out\textquotedblright{}.) Say it once more.
\end{tcolorbox}
\section[\texorpdfstring{\sk{संगृह्णाति} (saṅgṛhṇāti)}{संगृह्णाति (saṅgṛhṇāti)}]{\sk{संगृह्णाति} (saṅgṛhṇāti) — collects}

\label{lesson:SA-C177-sangrhnati}



\begin{quote}
Before the new one: say the Sanskrit for conceals, then the Sanskrit for distributes.
\end{quote}

\subsection*{You'll want to know: \sk{संगृह्णाति}}
\textbf{\sk{संगृह्णाति}} — \emph{saṅgṛhṇāti} — \textquotedblleft{}collects\textquotedblright{}.

\begin{grammarlens}[title={What you've built}]
One more word for this chapter.
\end{grammarlens}

\subsection*{Your turn}
\begin{itemize}
\raggedright
  \item \emph{Say it:} \emph{saṅgṛhṇāti}
  \item \emph{Say it:} \emph{saṅgṛhṇāti}, once more
  \item \emph{Say it:} say \emph{saṅgṛhṇāti}
  \item \emph{Recall:} say the Sanskrit for conceals, then the Sanskrit for distributes, then say \emph{saṅgṛhṇāti} again
\end{itemize}

\begin{tcolorbox}[breakable,colback=teal!4,colframe=teal!35!black,title={Before you move on}]
What does \sk{संगृह्णाति} mean? (\textquotedblleft{}Collects\textquotedblright{}.) Say it once more.
\end{tcolorbox}
\section[\texorpdfstring{\sk{विश्वसिति} (viśvasiti)}{विश्वसिति (viśvasiti)}]{\sk{विश्वसिति} (viśvasiti) — trusts}

\label{lesson:SA-C177-visvasiti}



\begin{quote}
Before the new one: say the Sanskrit for distributes, then the Sanskrit for collects.
\end{quote}

\subsection*{You'll want to know: \sk{विश्वसिति}}
\textbf{\sk{विश्वसिति}} — \emph{viśvasiti} — \textquotedblleft{}trusts\textquotedblright{}.

\begin{grammarlens}[title={What you've built}]
One more word for this chapter.
\end{grammarlens}

\subsection*{Your turn}
\begin{itemize}
\raggedright
  \item \emph{Say it:} \emph{viśvasiti}
  \item \emph{Say it:} \emph{viśvasiti}, once more
  \item \emph{Say it:} say \emph{viśvasiti}
  \item \emph{Recall:} say the Sanskrit for distributes, then the Sanskrit for collects, then say \emph{viśvasiti} again
\end{itemize}

\begin{tcolorbox}[breakable,colback=teal!4,colframe=teal!35!black,title={Before you move on}]
What does \sk{विश्वसिति} mean? (\textquotedblleft{}Trusts\textquotedblright{}.) Say it once more.
\end{tcolorbox}
\section[\texorpdfstring{\sk{लज्जते} (lajjate)}{लज्जते (lajjate)}]{\sk{लज्जते} (lajjate) — is ashamed, is shy}

\label{lesson:SA-C177-lajjate}



\begin{quote}
Before the new one: say the Sanskrit for collects, then the Sanskrit for trusts.
\end{quote}

\subsection*{You'll want to know: \sk{लज्जते}}
\textbf{\sk{लज्जते}} — \emph{lajjate} — \textquotedblleft{}is ashamed, is shy\textquotedblright{}.

\begin{grammarlens}[title={What you've built}]
One more word for this chapter.
\end{grammarlens}

\subsection*{Your turn}
\begin{itemize}
\raggedright
  \item \emph{Say it:} \emph{lajjate}
  \item \emph{Say it:} \emph{lajjate}, once more
  \item \emph{Say it:} say \emph{lajjate}
  \item \emph{Recall:} say the Sanskrit for collects, then the Sanskrit for trusts, then say \emph{lajjate} again
\end{itemize}

\begin{tcolorbox}[breakable,colback=teal!4,colframe=teal!35!black,title={Before you move on}]
What does \sk{लज्जते} mean? (\textquotedblleft{}Is ashamed, is shy\textquotedblright{}.) Say it once more.
\end{tcolorbox}
